\section{The Style Section: JavaScript/TypeScript}
\label{sec:js-compiler}

JavaScript gets special focus in Emi, because nearly everything eventually
has to be consumed by a JS environment. The JS output is real, runnable
JavaScript --- no transpile step required --- while \texttt{--tags
typescript} produces high-coverage \texttt{.ts} with, Emi's docs claim,
some of the highest-quality TypeScript output among competing generators.
\texttt{--tags react} layers on TanStack React Query hooks plus an
internal \texttt{useSse} and \texttt{useWebSocketX}; \texttt{--tags
nestjs} emits decorators that make request, headers and query params
portable straight into Nest.js handlers with no extra code.

\subsection{Auto-init: every field guaranteed to exist}
When an action returns a class instance, \emph{every} non-nullable field
--- including nested objects, arrays, and child DTOs --- is guaranteed to
exist. Writing this by hand is a hassle and easy to get subtly wrong;
Emi generates it from the definition alone:
\begin{lstlisting}
fields:
  - name: object1
    type: object
    fields:
      - name: contacts
        type: array
        fields:
          - name: email
            type: string
            default: emi-compiler@emi-compiler.com
          - name: phone
            type: string?
\end{lstlisting}
produces a class whose \texttt{object1} field is \emph{always} an
instance of a nested \texttt{AutoInitClassDto.Object1} class, never
\texttt{undefined} --- the code stops sprinkling \texttt{?.}, \texttt{??},
and \texttt{typeof x === 'string'} defensively, because the class already
did the guarding.

\subsection{Ergonomics every generated class ships with}
\begin{itemize}[leftmargin=1.1em,itemsep=2pt,topsep=2pt]
  \item \texttt{from()} --- full object constructor; \texttt{with()} ---
        partial, deeply-typed via \texttt{PartialDeep}.
  \item \texttt{copyWith()}, \texttt{clone()}, a correct \texttt{toJSON()}
        that sees the private fields (plain \texttt{JSON.stringify}
        otherwise can't --- the fields are genuinely private).
  \item A static \texttt{Fields} map (e.g.\ \texttt{Fields.date === "date"})
        for safe, refactor-proof path access instead of a hardcoded string
        literal.
\end{itemize}

\subsection{The Headers class}
HTTP request/response headers are a key/value definition; Emi generates a
class extending the standard \texttt{Headers} with typed
\texttt{getX()}/\texttt{setX()} accessors that coerce values on the way
out:
\begin{lstlisting}
- {name: accept-language, type: string}
- {name: cache-time,      type: int64}
\end{lstlisting}
\begin{lstlisting}
export class SampleHeaders extends Headers {
  getAcceptLanguage() { return this.#getTyped('accept-language', 'string|null'); }
  setCacheTime(value: number|null) {
    if (value !== null) this.set('cache-time', value);
    return this;
  }
}
\end{lstlisting}
\texttt{.get()} on the plain \texttt{Headers} base still returns the raw
string --- the typed accessors are additive, not a replacement.

\begin{pullquote}
Header/query-param typing at this level --- including \emph{nested} query
params rendered as typed accessors --- is a feature Emi's own docs single
out as unique among the generators it's compared against.
\end{pullquote}
